% Options for packages loaded elsewhere
\PassOptionsToPackage{unicode}{hyperref}
\PassOptionsToPackage{hyphens}{url}
\documentclass[
  ignorenonframetext,
]{beamer}
\newif\ifbibliography
\usepackage{pgfpages}
\setbeamertemplate{caption}[numbered]
\setbeamertemplate{caption label separator}{: }
\setbeamercolor{caption name}{fg=normal text.fg}
\beamertemplatenavigationsymbolsempty
% remove section numbering
\setbeamertemplate{part page}{
  \centering
  \begin{beamercolorbox}[sep=16pt,center]{part title}
    \usebeamerfont{part title}\insertpart\par
  \end{beamercolorbox}
}
\setbeamertemplate{section page}{
  \centering
  \begin{beamercolorbox}[sep=12pt,center]{section title}
    \usebeamerfont{section title}\insertsection\par
  \end{beamercolorbox}
}
\setbeamertemplate{subsection page}{
  \centering
  \begin{beamercolorbox}[sep=8pt,center]{subsection title}
    \usebeamerfont{subsection title}\insertsubsection\par
  \end{beamercolorbox}
}
% Prevent slide breaks in the middle of a paragraph
\widowpenalties 1 10000
\raggedbottom
\AtBeginPart{
  \frame{\partpage}
}
\AtBeginSection{
  \ifbibliography
  \else
    \frame{\sectionpage}
  \fi
}
\AtBeginSubsection{
  \frame{\subsectionpage}
}
\usepackage{iftex}
\ifPDFTeX
  \usepackage[T1]{fontenc}
  \usepackage[utf8]{inputenc}
  \usepackage{textcomp} % provide euro and other symbols
\else % if luatex or xetex
  \usepackage{unicode-math} % this also loads fontspec
  \defaultfontfeatures{Scale=MatchLowercase}
  \defaultfontfeatures[\rmfamily]{Ligatures=TeX,Scale=1}
\fi
\usepackage{lmodern}
\ifPDFTeX\else
  % xetex/luatex font selection
\fi
% Use upquote if available, for straight quotes in verbatim environments
\IfFileExists{upquote.sty}{\usepackage{upquote}}{}
\IfFileExists{microtype.sty}{% use microtype if available
  \usepackage[]{microtype}
  \UseMicrotypeSet[protrusion]{basicmath} % disable protrusion for tt fonts
}{}
\makeatletter
\@ifundefined{KOMAClassName}{% if non-KOMA class
  \IfFileExists{parskip.sty}{%
    \usepackage{parskip}
  }{% else
    \setlength{\parindent}{0pt}
    \setlength{\parskip}{6pt plus 2pt minus 1pt}}
}{% if KOMA class
  \KOMAoptions{parskip=half}}
\makeatother
\setlength{\emergencystretch}{3em} % prevent overfull lines
\providecommand{\tightlist}{%
  \setlength{\itemsep}{0pt}\setlength{\parskip}{0pt}}
\usepackage{bookmark}
\IfFileExists{xurl.sty}{\usepackage{xurl}}{} % add URL line breaks if available
\urlstyle{same}
\hypersetup{
  hidelinks,
  pdfcreator={LaTeX via pandoc}}

\author{\texorpdfstring{}{}}
\date{}

\begin{document}

\begin{frame}{The Atlantic Crucible: American Bimetallism and the
Architecture of Republican Money}
\protect\phantomsection\label{the-atlantic-crucible-american-bimetallism-and-the-architecture-of-republican-money}
In 1793---the same year the British state declared war on revolutionary
France, the same year radical publishers were arrested under Pitt's
sedition laws---William Blake etched \emph{America A Prophecy} on
eighteen copper plates in his Lambeth workshop. Blake was no isolated
mystic: David Erdman documents that he moved through the radical circles
of 1790s London, frequenting the weekly dinners at Joseph Johnson's
bookshop in St Paul's Churchyard alongside Mary Wollstonecraft, William
Godwin, and Thomas Paine {[}@erdman1954prophet{]}. When the Pitt
government issued a warrant for Paine's arrest for seditious libel after
\emph{Rights of Man} (1791--92), Blake personally warned Paine, enabling
his escape to France (Erdman p.~200). The poem that emerged from this
milieu presents the American Revolution not as a political event but as
a cosmic conflagration: the unchaining of Orc, the spirit of
revolutionary energy, against the tyranny of Urizen's law. But
\emph{America A Prophecy} is not merely a poem about political
revolution; it is, when read alongside the concurrent monetary crises of
both Britain and the newly constituted United States, a prophetic
meditation on the fundamental relationship between monetary sovereignty,
human freedom, and the architecture of value itself.

\begin{block}{The Coinage Act of 1792: Hamilton's Republican Experiment}
\protect\phantomsection\label{the-coinage-act-of-1792-hamiltons-republican-experiment}
On 2 April 1792---one year before Blake etched
\emph{America}---President George Washington signed the Coinage Act of
1792, establishing the United States Mint in Philadelphia and codifying
a bimetallic standard with a gold-to-silver ratio of \textbf{15:1}
{[}@laughlin1898bimetallism{]}. This ratio, proposed by Treasury
Secretary Alexander Hamilton in his \emph{Report on the Establishment of
a Mint} (1791), was itself an act of deliberate divergence from Newton's
British ratio of 1:15.21 {[}@hamilton1791report{]}. Hamilton explicitly
calculated his ratio by surveying average European market conditions,
aiming for a standard that would attract both metals into domestic
circulation rather than suffer the asymmetric drainage that plagued
Britain {[}@sylla2011hamiltons{]}. The silver dollar was defined at
\textbf{371.25 grains of pure silver}; the gold eagle at \textbf{247.5
grains of pure gold}. As Richard Sylla and David Cowen demonstrate,
Hamilton's financial architecture---encompassing the Mint, the First
Bank, and the systematic funding of revolutionary war debt---was the
most ambitious experiment in republican finance the modern world had yet
attempted {[}@sylla2011hamiltons{]}.

This legislative act was revolutionary in its monetary implications, yet
it carried a structural flaw that would destabilize the American economy
for a century. Where Newton's Mint had served the Crown's imperial
prerogative---calibrating a 15.21:1 ratio that systemically overvalued
gold and drained England of silver---Hamilton designed his 15:1 ratio
for a republican experiment in dual-metal equilibrium. However, as
Angela Redish demonstrates in \emph{Bimetallism: An Economic and
Historical Analysis}, Hamilton's flatter ratio sat blindly ``on the
slope of a declining value of silver relative to gold,'' driven by an
oversupply from Latin American mines {[}@redish2000bimetallism{]}.
Consequently, Hamilton inadvertently overvalued silver. This triggered
Gresham's Law in reverse: gold coins were melted and exported to Europe
where they commanded a higher premium, leaving the United States on a
\emph{de facto} silver standard by 1834. As Lawrence Officer calculates,
the U.S. Mint earned substantial seigniorage profit from the operation
of the bimetallic system despite its structural instability, suggesting
the arrangement served fiscal as well as monetary purposes
{[}@officer2001bimetallism{]}. Milton Friedman and Anna Schwartz later
analyzed this initial miscalibration as guaranteeing the monetary
volatility that would define 19th-century America
{[}@friedman1963monetary{]}. Curtis Nettels' study of the early national
economy confirms that Hamilton's system, however flawed in its ratio,
succeeded in establishing the institutional architecture---Mint,
Treasury, funded debt---upon which American monetary sovereignty would
depend {[}@nettels1962emergence{]}.

Blake's Orc, breaking free from his chains on the Atlantean mountains,
embodies precisely this revolutionary energy. The Preludium to
\emph{America} (Plates 1--4) depicts Orc's liberation as an elemental
eruption: the ``hairy youth'' bursts his chains in a subterranean
cavern, his body burning with the fire of revolutionary desire (Erdman
pp.~51--52) {[}@erdman1988complete{]}. When we read this mythological
unchaining alongside the simultaneous establishment of a new monetary
sovereign---a nation explicitly rejecting the Bank of England's
monopoly, the Crown's prerogative of issuance, and the Newtonian
metrological regime---Orc's fire becomes legible as the liberation of
value itself from the calculus of empire.
\end{block}

\begin{block}{Hamilton vs.~Jefferson: The Structural Antinomy of
Republican Finance}
\protect\phantomsection\label{hamilton-vs.-jefferson-the-structural-antinomy-of-republican-finance}
The American experiment in bimetallism was, from its inception, riven by
an internal antinomy that maps with remarkable precision onto Blake's
mythological framework. Alexander Hamilton---the architect of the First
Bank of the United States (chartered 1791)---envisioned a centralized
financial architecture modeled on, though deliberately constrained
relative to, the Bank of England. Hamilton's bank, capitalized at \$10
million, was authorized to issue banknotes, manage federal debt, and
coordinate fiscal policy---functions that Thomas Jefferson viewed as a
direct assault on republican agrarian virtue
{[}@laughlin1898bimetallism; @wilentz2005rise{]}.

Jefferson's opposition to the Bank was not merely constitutional (though
he argued strenuously that the Constitution granted no power to charter
a federal bank); it was ontological. Jefferson insisted that the
productive labor of the farmer---the physical working of the
earth---constituted the only legitimate foundation of national wealth.
Paper credit, fractional reserve banking, and the abstractions of urban
finance were, in Jefferson's view, parasitical fictions that enriched
speculators at the expense of the producing classes. As Sean Wilentz
documents in \emph{The Rise of American Democracy}, this
Hamilton-Jefferson divide defined the fundamental fault line of American
political economy for the next century, structuring every subsequent
conflict over banking, tariffs, and monetary standards
{[}@wilentz2005rise{]}. This position resonates with Blake's
condemnation of ``allegoric riches'' (\emph{Jerusalem}, Plate 27; Erdman
p.~170) {[}@erdman1988complete; @sklar2011blake{]}: both Blake and
Jefferson recognized that the architecture of money creation determined
the distribution of real power and real suffering.

Yet Jefferson's agrarian alternative was itself caught in the
contradictions of commodity production. The Jeffersonian farmer remained
dependent on metallic currency---on the very gold and silver whose
statutory ratios, as Gresham's Law demonstrated, were subject to
international arbitrage beyond any single nation's control. Blake's
prophetic critique extends further than either Hamilton's centralized
finance or Jefferson's agrarian metallist position: for Blake, the
fundamental error is the reduction of value to \emph{any} quantitative
metric, whether gold, silver, or paper. ``Can Wisdom be put in a silver
rod? / Or Love in a golden bowl?'' (\emph{The Book of Thel}, 1789)
{[}@blake1789thel{]}---this question applies with equal force to
Hamilton's Treasury notes and Jefferson's silver dollars.
\end{block}
\end{frame}

\end{document}
